% Capítulo 4: Materiales y métodos.
\label{cap:metodos}

Este capítulo describe los materiales y el procedimiento del estudio
con el detalle necesario para reproducirlo. Se presentan primero los
conjuntos de datos y los métodos clásicos de selección que forman la
referencia, junto con la construcción del grafo de características de
relevancia y redundancia. A continuación se detalla el método cuántico
basado en el bloqueo de Rydberg, con su codificación física y su
contrato de evaluación. El capítulo concluye con el diseño
experimental, que fija las fases del trabajo, el control de fuga de
información, el protocolo de modelado y los contrastes estadísticos.

\section{Conjuntos de datos}
\label{sec:datos}

Se seleccionaron cuatro conjuntos de datos de clasificación
supervisada, elegidos para cubrir regímenes distintos de tamaño
muestral, dimensionalidad y tipo de variables
(Tabla~\ref{tab:datasets}). \textit{Breast Cancer Wisconsin}
(diagnóstico)~\cite{street1993} es un problema binario clínico con
fuerte redundancia entre medidas morfológicas de una misma familia.
\textit{Customer Churn} es un problema binario de gran volumen
muestral con variables numéricas y categóricas~\cite{churn_kaggle}.
\textit{Madelon}~\cite{guyon2004} es un problema sintético de alta
dimensionalidad diseñado para el desafío de selección de
características de NIPS~2003: de sus 500 variables solo 20 son
informativas o redundantes por construcción y las 480 restantes son
distractores. Por último, \textit{Olive Oil}~\cite{forina1983}
contiene la composición en ácidos grasos de 572 aceites de oliva
italianos, etiquetados a dos niveles geográficos. Esta diversidad
extiende deliberadamente el banco de pruebas de la propuesta original
del método cuántico, evaluado sobre tres problemas binarios de
dimensionalidad moderada~\cite{orquin2026}, hacia escenarios
multiclase, desbalanceados y de alta dimensionalidad con verdad
conocida.

La elección de los conjuntos no responde solo a variedad descriptiva.
Cada uno ejerce presión sobre un componente distinto del proceso de
selección. \textit{Breast Cancer Wisconsin} concentra señal fuerte en
familias de tamaño celular altamente redundantes; \textit{Customer
Churn} combina un gran tamaño muestral con variables continuas
dominantes y categóricas nominales; \textit{Madelon} fuerza un régimen
de alta dimensionalidad con señal no univariante; y \textit{Olive Oil}
permite observar un espacio químico pequeño, composicional y multiclase.
Esta lectura por regímenes se conserva durante todo el diseño: las
fases clásicas no solo preparan los datos, sino que anticipan qué parte
del método QFS debería verse tensionada en cada caso.

\begin{table}[!htbp]
  \centering
  \caption[Conjuntos de datos del estudio]{Conjuntos de datos del estudio. El recuento de predictoras
    excluye identificadores estructurales y, en \textit{Olive Oil},
    las etiquetas alternativas y su proxy (véase la
    sección~\ref{sec:oliveoil}).}
  \label{tab:datasets}
  \begin{tabular}{lrrll}
    \toprule
    Dataset & Casos & Variables predictoras & Target & Clases \\
    \midrule
    Breast Cancer Wisconsin & 569 & 30 & diagnóstico & 2 \\
    Customer Churn & 440\,832 & 10 & abandono & 2 \\
    Madelon & 2\,000 & 500 & sintético & 2 \\
    Olive Oil & 572 & 8 & región de origen & 3 / 9 \\
    \bottomrule
  \end{tabular}
\end{table}

\subsection{Decisiones específicas por conjunto de datos}
\label{sec:oliveoil}

Dos conjuntos exigen decisiones de tratamiento propias que se fijan
desde el particionado y se declaran aquí para no contaminar las
comparaciones posteriores: la separación de las etiquetas de
\textit{Olive Oil} y la codificación de las variables categóricas de
\textit{Customer Churn}.

El análisis exploratorio mostró que \textit{Olive Oil} contiene dos
etiquetas supervisadas jerárquicas, la macro-región (3 clases) y la
región (9 clases), además de una variable (\texttt{palmitic} en la
versión utilizada) que actúa como sustituto (\textit{proxy}) casi
determinista de la región. Tratar el conjunto como un único problema
habría contaminado cualquier comparación, porque la etiqueta
alternativa y el proxy quedarían dentro del espacio de predictores.
Por ello, desde la fase de particionado el conjunto se desdobla en dos
formulaciones operativas no intercambiables:
\texttt{olive\_oil\_3class} (target macro-región; se excluyen la
etiqueta de 9 clases y el proxy) y \texttt{olive\_oil\_9class} (target
región; se excluyen la macro-región y el proxy). Todas las fases
posteriores trabajan con cinco formulaciones: los tres conjuntos
originales restantes más estas dos variantes, que nunca se promedian
entre sí. Tras estas exclusiones, cada formulación de \textit{Olive Oil}
trabaja con ocho predictoras de composición (las nueve variables de
ácidos grasos del conjunto crudo, una vez excluido el proxy
\texttt{palmitic}). El recuento de diez variables que aparece en el
análisis exploratorio incluye, además, la etiqueta jerárquica que cada
formulación no emplea como objetivo y que solo se descarta al fijarla; de
ahí que la tabla de señal liste diez variables y la de conjuntos de datos,
ocho.

\textit{Customer Churn} contiene tres predictoras categóricas nominales:
género, tipo de suscripción y duración del contrato. En los pipelines
posteriores al particionado se codifican mediante \textit{one-hot}
ajustado solo con entrenamiento, en lugar de imponer un orden artificial
mediante \textit{label encoding}. Esta decisión es metodológicamente
coherente con modelos lineales y con la comparación clásica, pero se
declara como divergencia respecto al \textit{label encoding} empleado en
la referencia QFS original~\cite{orquin2026}. Su efecto se comprueba después: genera grupos de
dummies mutuamente excluyentes que pueden competir por el presupuesto de
QFS, pero no se asume como causa del resultado sin comprobar el embebido
geométrico.

\section{Métodos clásicos de referencia}
\label{sec:clasicos}

Los métodos clásicos no se eligen como un muestrario genérico de
selectores, sino como el \emph{espejo} de la estructura del método
cuántico que se quiere evaluar. La función de coste de ese método
(sección~\ref{sec:grafo}) combina dos ingredientes, la relevancia de
cada variable con el target $I(x_i;y)$ y la redundancia entre pares
$I(x_i;x_j)$, y los resuelve de forma global. Por ello se selecciona al
menos un método clásico que aísle o combine cada ingrediente, más los
benchmarks de referencia del campo. El objetivo es que la comparación
posterior sea \emph{atribuible}: poder distinguir si el método cuántico
aporta por la relevancia, por el control de la redundancia o por
resolver el compromiso de forma global, en lugar de un resultado
agregado sin diagnóstico. El repertorio resultante consta de doce
selectores agrupados en seis roles:

\begin{itemize}
\item \textbf{Baseline mínimo}: la varianza (\texttt{variance}),
  filtro no supervisado que se conserva como suelo de comparación.
  Siguiendo el protocolo habitual de la selección no
  supervisada~\cite{solorio2020}, que estandariza las variables antes de
  aplicar los filtros, su puntuación opera sobre variables de varianza
  homogénea y apenas discrimina entre ellas; por eso se interpreta como
  una referencia inferior degenerada, no como un criterio de selección.
\item \textbf{Relevancia pura}: la puntuación F de ANOVA
  (\texttt{f\_classif}) y la información mutua
  (\texttt{mutual\_info})~\cite{cover2006}, que cuantifican la
  dependencia de cada variable con el target. La información mutua es,
  además, la cantidad exacta que el método cuántico codifica como
  relevancia.
\item \textbf{Redundancia pura}: la correlación mutua
  (\texttt{mutual\_correlation}) y la selección por similitud entre
  variables (\texttt{feature\_similarity})~\cite{solorio2020}, filtros
  no supervisados que eliminan variables redundantes sin mirar el
  target; representan el segundo ingrediente del método cuántico. La
  evaluación sistemática de Solorio-Fernández et al.~\cite{solorio2020}
  los sitúa entre los filtros no supervisados de menor rendimiento
  predictivo, por lo que su inclusión responde a su papel como
  \emph{espejo estructural} del término de redundancia y no a una
  expectativa de competitividad en exactitud.
\item \textbf{Relevancia y redundancia combinadas}: una aproximación
  voraz al criterio mRMR~\cite{peng2005} (\texttt{mrmr\_approx}) y la
  selección relevancia-redundancia (\texttt{rrfs})~\cite{solorio2020},
  los análogos clásicos directos del compromiso que el método cuántico
  resuelve de forma global.
\item \textbf{Benchmarks \textit{wrapper}}: Boruta~\cite{kursa2010},
  referencia \textit{gold-standard} de la selección, y la eliminación
  recursiva de variables (\texttt{rfe}). Boruta confirma el conjunto de
  todas las variables relevantes mediante un contraste frente a
  variables sombra, por lo que devuelve un subconjunto de tamaño
  determinado por los datos, no por un presupuesto $k$ prefijado.
\item \textbf{Métodos \textit{embedded}}: la regresión logística con
  penalización $\ell_1$, la importancia por impureza de \textit{random
  forest}~\cite{breiman2001} y la magnitud de pesos de una SVM lineal,
  que seleccionan en el contexto de un modelo.
\end{itemize}

Esta taxonomía permite leer los doce métodos como un sistema de
coordenadas. Los filtros de relevancia indican qué ocurre si solo se
atiende a $I_i$; los métodos de redundancia muestran el efecto de
penalizar similitud sin mirar el target; mRMR y RRFS aproximan el mismo
compromiso que QFS, pero de forma clásica; y los métodos embedded y
wrappers actúan como referencia de lo que un modelo puede recuperar por
otra vía. En el capítulo de resultados, QFS se situará en este plano
comparando tanto coordenadas agregadas (relevancia capturada y
redundancia interna) como subconjuntos por nombre de variable.

Salvo Boruta, cada método produce un ranking del que se derivan
subconjuntos de tamaño $k$ según una escalera adaptada a cada dataset.
El $k$ de referencia para el modelado es el mayor $k \le 10$ que implica
una reducción real del espacio: 10 en los tres conjuntos grandes, y 5 en
las dos formulaciones de \textit{Olive Oil}, donde un presupuesto de 10
no reduciría nada porque solo hay 8 predictoras. Boruta se evalúa al
tamaño natural de su conjunto confirmado. La cota superior $k \le 10$ es
coherente con la evidencia de la referencia cuántica, en la que el
rendimiento se estabiliza más allá de unas siete variables y la ventaja
del método se concentra en los subconjuntos
compactos~\cite{orquin2026}. La información mutua se estima
con la convención única descrita en la sección~\ref{sec:grafo}
(discretización en cinco intervalos uniformes), compartida por los
selectores basados en información mutua y por las matrices que alimentan
el método cuántico. Para acotar el coste, los métodos no lineales se
ajustan en \textit{Customer Churn} sobre una muestra estratificada de
5\,000 filas de entrenamiento cuya representatividad se verifica.

Además del ranking se miden tres propiedades de calidad. La
\emph{estabilidad} entre semillas (tres semillas independientes) se
cuantifica con el índice de Jaccard entre dos subconjuntos $A$ y $B$,
\begin{equation}
  \label{eq:jaccard}
  J(A,B) = \frac{|A \cap B|}{|A \cup B|},
\end{equation}
y con el índice de Kuncheva~\cite{kuncheva2007}, que para dos
subconjuntos de tamaño $k$ con $r$ variables en común, extraídos de un
espacio de $p$ variables, corrige el solape esperado por azar:
\begin{equation}
  \label{eq:kuncheva}
  K = \frac{r - k^2/p}{k - k^2/p}.
\end{equation}
A diferencia de Jaccard, $K$ vale cero cuando el solape coincide con el
esperado al azar y uno cuando es total. La \emph{separación frente
al azar} contrasta la hipótesis nula de que la puntuación de un selector
no es mayor que la que obtendría sobre un target sin relación con las
variables. Esa distribución nula se construye empíricamente permutando
el target 20 veces y recalculando las puntuaciones; una variable se
considera señal solo si su puntuación real supera el percentil 95 de la
distribución nula (resolución mínima del p-valor empírico
$1/21 \approx 0{,}048$). Este control se emplea como pantalla operativa
de separación frente al azar en la fase de selección; la evidencia
inferencial principal se obtiene después sobre las predicciones de prueba
mediante \textit{bootstrap} y permutaciones pareadas. La \emph{redundancia interna} de cada subconjunto
se mide como el cambio de la correlación absoluta media entre sus
variables respecto al espacio completo.

\section{Construcción del grafo de características}
\label{sec:grafo}

El puente entre la selección clásica y la cuántica es un grafo de
características construido exclusivamente con la partición de
entrenamiento. Para cada formulación se calculan el vector de
relevancias $I_i = I(x_i; y)$, la información mutua de cada variable
con el target, y la matriz de redundancias $R_{ij} = I(x_i; x_j)$, la
información mutua entre pares de variables~\cite{cover2006}. La
estimación de la información mutua replica la de la implementación de
referencia~\cite{orquin2026}: las variables continuas se discretizan en
cinco intervalos uniformes y la información mutua se estima sobre los
datos discretizados, de forma consistente para los selectores clásicos
basados en información mutua y para las matrices que alimentan el método
cuántico, de modo que ambos comparten exactamente la misma cantidad. Las
variables son los nodos, ponderados por $I_i$; las redundancias altas
definen las aristas. Sobre ese grafo, el compromiso
relevancia--redundancia se expresa como la función de coste
\begin{equation}
  \label{eq:qfs}
  Q(x;\alpha) \;=\; -\alpha \sum_{i} I_i\, x_i \;+\;
  (1-\alpha) \sum_{i<j} R_{ij}\, x_i x_j,
\end{equation}
donde $x_i \in \{0,1\}$ indica si la variable $i$ se selecciona y
$\alpha \in [0,1]$ regula el compromiso. Se trata de la instancia
QUBO descrita en la sección~\ref{sec:mt-qubo}; su materialización en
átomos neutros aproxima un problema MWIS físico al convertir
redundancias altas en proximidad y bloqueo. Es también la versión global
del criterio que \texttt{mrmr\_approx} aproxima de forma voraz.

El vector y la matriz de cada formulación se calculan durante la fase
de selección clásica. De este
modo, el método cuántico consume exactamente las mismas cantidades
informacionales que los métodos clásicos de referencia, y cualquier
diferencia de resultados será atribuible al procedimiento de
optimización y no a las entradas.

\section{Método cuántico basado en el bloqueo de Rydberg}
\label{sec:qfs}

El método QFS evaluado en este trabajo~\cite{orquin2026} resuelve la
ecuación~\ref{eq:qfs} sobre una matriz de átomos neutros, aplicando la
correspondencia descrita en el marco teórico
(secciones~\ref{sec:mt-adiabatica} y~\ref{sec:mt-rydberg}).

\subsection{Codificación física y protocolo}

Cada variable se representa con un átomo. La relevancia $I_i$ se
codifica en la amplitud del \textit{detuning} local aplicado a cada
átomo, y la redundancia $R_{ij}$ se traduce en distancias espaciales:
en la correspondencia física de referencia, las parejas muy
redundantes se colocan más cerca mediante una distancia base
$d_{ij}\propto R_{ij}^{-1/6}$ y, por tanto, quedan más expuestas al
bloqueo de Rydberg, que impide su excitación simultánea y actúa como
restricción dura de exclusión~\cite{orquin2026}. En este trabajo esa
correspondencia se conserva como núcleo del embebido, pero se evalúa
una variante operativa trazada que añade un término de separación
dependiente de la relevancia. En la ecuación se usan las cantidades
normalizadas $\widehat R_{ij}$ y $\widehat I_i$, con el término de
relevancia activo solo cuando $\beta>0$:
\begin{equation}
  \label{eq:qfs-distancia-beta}
  \tilde d_{ij} = \widehat R_{ij}^{-1/6} +
  \beta\,(1+\widehat I_i)(1+\widehat I_j),
  \qquad d_{ij}\in[d_{\min},d_{\max}].
\end{equation}
El parámetro $\beta$ no se interpreta como una nueva magnitud física,
sino como un grado de libertad experimental para estudiar si separar
pares de variables relevantes modifica las densidades de Rydberg y la
selección final. En la implementación, $I_i$ y $R_{ij}$ se normalizan a
$[0,1]$ con mínimos y máximos calculados en entrenamiento, de modo que
la potencia inversa actúa sobre la redundancia ya normalizada. El caso
$\widehat R_{ij}=0$ (sin redundancia estimada entre el par) se trata
aparte: a esa pareja se le asigna directamente la distancia máxima
$d_{\max}$, que representa separación máxima sin bloqueo, sin evaluar la
potencia inversa, que en ese punto divergiría; la expresión anterior
solo rige para los pares con redundancia no nula. La matriz cruda
$\tilde d_{ij}$ se reescala después al intervalo $[d_{\min},d_{\max}]$
que determina el radio de bloqueo y solo entonces se aplica escalado
multidimensional (MDS), reteniendo embebidos que satisfacen la
separación mínima entre átomos.

El sistema evoluciona adiabáticamente desde el estado fundamental hasta
el hamiltoniano final mediante un protocolo temporal de dos pasos sobre
los campos de control: la frecuencia de Rabi $\Omega(t)$, que mide la
intensidad del acoplamiento láser; el \textit{detuning} global
$\Delta(t)$, que desplaza por igual la energía de todos los átomos; y el
\textit{detuning} local, que aplica a cada átomo un sesgo proporcional a
su relevancia $I_i$~\cite{orquin2026}. La
Figura~\ref{fig:qfs-adiabatica} muestra estos tres perfiles a lo largo
de los 4~$\mu$s de evolución. La medición repetida produce
una distribución de cadenas de bits, y las densidades de Rydberg
medias sobre las configuraciones de menor energía proporcionan un
ranking de variables del que se derivan subconjuntos para cada
cardinalidad $k$. Conviene aclarar el papel de $\alpha$ en este
contexto: el recorrido de $\alpha$ en el QUBO puro produce una escalera
discreta de óptimos, pero esa escalera puede saltar presupuestos. Por
ello, el control exacto usado en este trabajo combina el barrido de
$Q^*(\alpha)$ con una restricción explícita de cardinalidad cuando debe
compararse a un $k$ fijado. La implementación de referencia sobre átomos
neutros~\cite{orquin2026} fija $\alpha$ en un valor representativo y
obtiene los subconjuntos por corte top-$k$ sobre la densidad de Rydberg
media. En este trabajo se evalúan ambos mecanismos: se ejecuta el
protocolo operativo del simulador analógico con $\alpha$ fija, y se
reproduce además el recorrido de $\alpha$ propuesto por
Mücke~\cite{mucke2023} resolviendo $Q^*(\alpha)$ por enumeración
exhaustiva o programación entera dentro de la envolvente operativa
(sección~\ref{sec:contrato}). El segundo procedimiento no compite con
el primero: actúa como referencia exacta del criterio QUBO sobre la
misma función de coste, lo que permite separar la calidad del
optimizador analógico de la calidad del propio criterio. El segundo
grado de libertad es el parámetro $\beta$ introducido en la
ecuación~\ref{eq:qfs-distancia-beta}, que se barre por validación en
$\{0, 0{,}1, 0{,}2, \ldots, 1\}$. Esta extensión se declara separadamente
del mapeo base del artículo para evitar una lectura de reproducción
literal: el objetivo no es cambiar el criterio QUBO, sino comprobar si
la geometría operacional del simulador analógico altera el
\textit{readout}. La dinámica se simula con el módulo de simulación
analógica hamiltoniana de Amazon Braket en su simulador local.

\begin{figure}[!htbp]
  \centering
  \includegraphics[width=0.86\textwidth]{ev5_evolucion_adiabatica.png}
  \caption[Perfiles temporales de los campos de control]{Perfiles temporales de los campos de control durante los
    4~$\mu$s de evolución adiabática simulada. El eje horizontal es el
    tiempo, en microsegundos; las curvas muestran la frecuencia de Rabi
    $\Omega(t)$, el \textit{detuning} global $\Delta(t)$ y el
    \textit{detuning} local que codifica la relevancia de cada variable.
    Su forma escalonada en dos pasos es la que lleva el sistema desde el
    estado fundamental inicial hasta la configuración de mínima energía
    que codifica el subconjunto seleccionado.}
  \label{fig:qfs-adiabatica}
\end{figure}

\subsection{Envolvente operativa y contrato de evaluación}
\label{sec:contrato}

La simulación clásica de la dinámica analógica limita el número de
átomos tratables, y la implementación de referencia se ha validado con
problemas de hasta una veintena de variables~\cite{orquin2026}. La
envolvente operativa queda por tanto fijada en torno a veinte
variables: las formulaciones que caben en ella (\textit{Olive Oil},
con 8 variables, y \textit{Customer Churn}, con 10 predictoras crudas
que se materializan como 15 columnas tras la codificación
\textit{one-hot} ajustada en entrenamiento) entran directamente; las
que la exceden (\textit{Breast Cancer}, 30;
\textit{Madelon}, 500) requieren una preselección clásica previa que
reduce el espacio a unas 20 variables antes de la fase cuántica. Esa
preselección emplea \texttt{mrmr\_approx}, el análogo voraz del mismo
criterio de relevancia y redundancia que usa el método cuántico, para no
sesgar la entrada hacia un criterio distinto del que se evalúa; se ajusta
solo con entrenamiento y se documenta como parte del método híbrido
resultante. Esta elección evita introducir un criterio ajeno, pero no
elimina una dependencia importante: si la preselección descarta variables
verdaderamente relevantes (algo posible en \textit{Madelon}, donde la
señal reside en interacciones que un criterio par a par no captura), el
método cuántico no puede recuperarlas y hereda ese sesgo desde el inicio.
El alcance de este efecto se discute con los resultados de \textit{Madelon}
y se recoge entre las limitaciones; una mitigación natural sería
preseleccionar con la unión de varios criterios o con uno sensible a
interacciones de orden superior. Los parámetros de simulación se heredan
de la implementación de referencia~\cite{orquin2026}: una evolución
adiabática de 4~$\mu$s (tiempo total de evolución, elegido dentro del
régimen operativo considerado en dicha implementación), la retención del
10\,\% de configuraciones de menor
energía y 100
inicializaciones independientes del escalado multidimensional. El número
de mediciones por configuración se eleva de las 1\,000 de la referencia
a 10\,000 para reducir la varianza muestral del histograma de
\textit{bitstrings}.

Hay una condición geométrica del protocolo de átomos neutros. Tras el
embebido por MDS, las posiciones se reescalan para que la distancia
mínima entre dos átomos sea una fracción del radio de bloqueo $R_b$,
fijada en $1/\sqrt{2}$ como en la implementación de
referencia~\cite{orquin2026}; esa fracción determina cuántos pares
quedan dentro del volumen de bloqueo y, por tanto, cuántas exclusiones
impone la física. No siempre existe un embebido bidimensional que
respete esa separación, porque la matriz de redundancia no tiene por qué
ser realizable como grafo de discos unitarios~\cite{nguyen2023,
bermot2025}; en este banco, sin embargo, las cinco formulaciones
admiten un embebido válido a $1/\sqrt{2}$, que se conserva sin
relajación. La calidad del embebido se verifica después con el error MDS y
con la comparación frente al oráculo exacto, de modo que cualquier
deterioro de QFS no se atribuye a la geometría sin contrastarlo
explícitamente.

Las 100 inicializaciones del escalado multidimensional se realizan como
100 ejecuciones con semillas independientes. Para que el parámetro
interno \texttt{n\_init} de scikit-learn no multiplique de forma
implícita el número de arranques ni dependa de la versión de la
biblioteca, se fija \texttt{n\_init=1} dentro de ese bucle de 100
semillas, de modo que cada semilla equivale a un arranque independiente,
en correspondencia exacta con las 100 corridas del protocolo de
referencia. El error de embebido resultante queda trazado para el
análisis de resultados.

El tamaño de la envolvente habilita además un control de calidad
específico: con una veintena de variables, el óptimo exacto de la
ecuación~\ref{eq:qfs} es computable por enumeración exhaustiva, de
modo que puede medirse la distancia entre el subconjunto devuelto por
la evolución adiabática y el mínimo de la función de coste. Esta
comprobación, alineada con la validación mediante optimización exacta
empleada en la propuesta original~\cite{orquin2026}, separa dos
preguntas que la comparación final no debe confundir: si el método
cuántico optimiza bien su criterio y si ese criterio produce
subconjuntos predictivos. Su evidencia se materializa en el control de
calidad de la Tabla~\ref{tab:qfs-control} y en la descomposición
criterio--optimizador de la Tabla~\ref{tab:atribucion}
(sección~\ref{sec:diagnostico-qfs}).

El contrato de evaluación hereda íntegramente el diseño del bloque
clásico: las entradas del método son el grafo de características de la
sección~\ref{sec:grafo}. La configuración operativa sobre átomos neutros, QFS-NA, fija
$\alpha=0{,}5$ y obtiene la cardinalidad mediante un corte top-$k$ sobre
la densidad de Rydberg, en lugar de recorrer $\alpha$. Ese recorrido,
que en el QUBO exacto produce una escalera discreta de presupuestos
observados (sección~\ref{sec:mt-qubo}), se reserva al oráculo exacto
como control del criterio y no se presenta como salida del simulador
analógico. Cuando la escalera salta el presupuesto de referencia, el
oráculo se calcula con la restricción explícita $\sum_i x_i=k$. Sobre
la configuración operativa se barre el parámetro $\beta$ y se elige su
valor exclusivamente con la partición de validación, usando
\textit{random forest} como evaluador fijo para no convertir el bloque
cuántico en otro barrido de modelos; el conjunto de prueba no interviene
en ninguna elección. Una vez fijados los subconjuntos, el rendimiento se
evalúa con el protocolo de la sección~\ref{sec:protocolo} (mismos
modelos, macro-F1, intervalos \textit{bootstrap}, comparaciones pareadas
y contrastes de permutación) sobre los presupuestos de variables
fijados en la sección~\ref{sec:clasicos} ($k=10$ en los conjuntos
grandes y $k=5$ en \textit{Olive Oil}).

\section{Diseño experimental}
\label{sec:diseno}

Sobre todo el diseño se aplica una regla transversal, cuyo fundamento
se ha presentado en la sección~\ref{sec:mt-validacion}: ningún
hallazgo del bloque se acepta sin contraste estadístico; toda batería
de contrastes univariantes se corrige por multiplicidad; todo
contraste se acompaña de su tamaño de efecto; y, donde la distribución
nula no es conocida, se construye empíricamente por permutación. Las
subsecciones siguientes concretan esa regla en cada eslabón de la
cadena experimental.

El diseño experimental se organiza como una cadena de diez fases. Las
siete primeras construyen una referencia clásica contrastada, las dos
siguientes ejecutan, evalúan e integran el método cuántico contra esa
referencia, y la última fase sintetiza el recorrido sin reentrenar
modelos ni reejecutar QFS\@. Cada fase produce salidas verificables que se
encadenan con la fase siguiente: particiones, subconjuntos
seleccionados, predicciones, contrastes estadísticos, diagnósticos de
calidad, errores de embebido y figuras de cierre.

\begin{enumerate}
\item \textbf{Análisis exploratorio del dato crudo}: estructura,
  calidad, distribución del target, normalidad, asociación
  variable--target con corrección por contrastes múltiples, magnitud
  del efecto, redundancia, dimensionalidad y riesgos de
  \textit{leakage} potencial.
\item \textbf{Preprocesado estructural}: normalización de nombres,
  exclusión de identificadores, codificación determinista del target y
  caracterización de valores atípicos, sin imputación, escalado ni
  codificación global de predictoras.
\item \textbf{Comprobación posterior al preprocesado}: verificación de que el
  preprocesado no alteró distribuciones, target, asociación
  variable--target ni estructura de redundancia.
\item \textbf{Particionado y verificación de las particiones}: creación de las
  particiones de entrenamiento (70\,\%), validación (15\,\%) y prueba
  (15\,\%), estratificadas por target, y verificación de solapes,
  conservación del target, \textit{drift} univariante,
  representatividad multivariante, \textit{leakage} y separabilidad
  adversarial.
\item \textbf{Selección clásica de características}: ejecución y
  caracterización de los doce selectores bajo contrato
  \textit{train-only}, incluida la construcción del grafo de
  características de la sección~\ref{sec:grafo}.
\item \textbf{Modelado y evaluación estadística}: comparación de los
  subconjuntos seleccionados frente al espacio completo con cuatro
  modelos y contrastes sobre el conjunto de prueba.
\item \textbf{Comparación final y mapa de evidencia}: integración de
  los resultados de las fases anteriores, veredicto por dataset y
  preparación de la referencia para la fase cuántica.
\item \textbf{Ejecución QFS}: preselección híbrida cuando la envolvente
  operativa lo exige, barrido de $\beta$, simulación analógica
  hamiltoniana, embebido MDS con 100 inicializaciones
  independientes, adaptación geométrica trazada y control frente al
  óptimo exacto del QUBO.
\item \textbf{Evaluación e integración cuántica}: evaluación de los
  subconjuntos QFS, modelado con el mismo protocolo de la referencia
  clásica, comparación pareada frente al baseline e integración en la
  tabla final.
\item \textbf{Síntesis visual y verificación final}: lectura conjunta de
  los artefactos ya generados para construir el diagnóstico
  multidimensional régimen--mecanismo--resultado, incluyendo posición de
  QFS frente a los doce métodos, trayectorias en $k$, escalera
  $\alpha$, efecto de $\beta$, error MDS, atribuciones SHAP
  (\textit{SHapley Additive exPlanations}), métricas y consistencia.
\end{enumerate}

\subsection{Preprocesado y comprobación de datos}

El preprocesado se diseña bajo un criterio de intervención mínima y
sin ajuste global: normalización de nombres, exclusión de
identificadores estructurales y codificación determinista del target,
sin imputación, winsorización, escalado ni codificación de predictoras
sobre el conjunto completo; esas transformaciones se posponen a los
pipelines posteriores al particionado. Cada decisión se comprueba en la
fase siguiente comparando el dataset procesado con el crudo en
dimensiones, distribución del target, distribuciones numéricas, orden
de asociación variable--target y estructura de redundancia. Este doble
paso garantiza que los problemas llegan al particionado sin
alteraciones que pudieran favorecer a ningún método de selección.

\subsection{Control de fuga de información}

Todo ajuste dependiente de los datos se realiza exclusivamente con la
partición de entrenamiento. En el preprocesado, las predictoras
categóricas se conservan sin codificar hasta después del particionado,
de modo que los codificadores \textit{one-hot} y los escaladores se
ajustan con entrenamiento y se aplican después a validación y prueba.
En la selección de características, cada selector se ajusta con
$X_{train}$ e $y_{train}$; validación y prueba solo reciben las
columnas ya seleccionadas. En el modelado, la elección de
configuraciones se realiza con validación y el conjunto de prueba se
utiliza una única vez, sobre candidatos fijados de antemano. Además,
se excluyen del espacio de predictores los identificadores, el índice
de trazabilidad de las particiones y, en \textit{Olive Oil}, las
etiquetas alternativas y su proxy.

\subsection{Validación estructural y de particiones}

La verificación de las particiones combina comprobaciones exactas y
diagnósticos estadísticos. Las exactas verifican ausencia de solapes
de filas entre particiones y conservación de todas las clases. Las
estadísticas comparan
entrenamiento frente a validación y prueba mediante: el contraste
$\chi^2$ y la diferencia máxima de proporciones para el target; los
estadísticos de Kolmogorov--Smirnov, la distancia de Wasserstein
estandarizada y el índice de estabilidad poblacional (PSI) para las
variables numéricas; y el contraste $\chi^2$ con distancia de
proporciones y PSI categórico para las variables discretas. Como
diagnóstico global de representatividad se emplea una validación
adversarial: un clasificador que intenta distinguir filas de
entrenamiento y prueba, cuyo AUC (área bajo la curva ROC) se interpreta como no accionable si
queda cerca de 0{,}5 bajo el umbral operativo fijado para la verificación.
Como pantalla de \textit{leakage} se calcula,
para cada variable, el AUC univariante absoluto y la información mutua
normalizada con el target, marcando para revisión los valores
superiores a 0{,}99. Los resultados de esta verificación se sintetizan en la
Tabla~\ref{tab:particiones} de la sección~\ref{sec:res-base}.

\subsection{Protocolo de modelado y contrastes estadísticos}
\label{sec:protocolo}

La fase de modelado cruza, por dataset, el espacio completo de
variables (\textit{baseline}) y los doce subconjuntos clásicos de
la fase de selección con cuatro modelos de familias distintas:
regresión logística, SVM lineal, \textit{random
forest}~\cite{breiman2001} y \textit{XGBoost}~\cite{chen2016}, todos
con hiperparámetros fijos y pesos de clase balanceados. No se realiza
búsqueda de hiperparámetros: el objeto de estudio es el espacio de
variables, y el baseline comparte modelo e hiperparámetros con cada
subconjunto, de modo que las diferencias observadas son atribuibles a
las variables. La inclusión de \textit{XGBoost} responde a un criterio
explícito de comparabilidad con la evaluación original del método
cuántico, que emplea ese mismo clasificador~\cite{orquin2026}, mientras
que las tres familias adicionales extienden la cobertura a modelos
lineales y de conjunto bajo el mismo protocolo. El método cuántico se
evaluará con estos mismos modelos (sección~\ref{sec:contrato}), de modo
que la comparación clásico--cuántico de este trabajo queda
internamente homogénea sin renunciar al puente con las cifras
preliminares de la referencia.

La métrica primaria es el macro-F1, que pondera por igual todas las
clases, con la exactitud balanceada como criterio de desempate. Esta
elección se justifica porque el banco de este trabajo incluye problemas
multiclase y desbalanceados en los que el AUC binario empleado en la
evaluación original del método cuántico~\cite{orquin2026} no es
aplicable de forma directa. En los dos conjuntos binarios del banco
(\textit{Breast Cancer Wisconsin} y \textit{Customer Churn}) se
registra adicionalmente el AUC-ROC como métrica de contexto frente a
las cifras preliminares de la referencia, pero no se utiliza como criterio primario de
decisión, que se mantiene homogéneo entre todos los datasets.
Por dataset pasan a la evaluación final el mejor baseline y los dos
mejores subconjuntos en validación. Como evidencia de explicabilidad,
para cada configuración evaluada se calculan los valores
SHAP~\cite{lundberg2017} con el explicador nativo de cada familia de
modelo (\textit{TreeExplainer} para los bosques aleatorios y
\textit{LinearExplainer} para la regresión logística y la SVM lineal),
estimados sobre la partición de validación con un fondo tomado del
entrenamiento, junto con las matrices de confusión de prueba; estos
resultados permiten comprobar qué variables sostienen cada subconjunto
seleccionado.

Sobre las predicciones de prueba se calculan: intervalos de confianza
del 95\,\% por \textit{bootstrap}~\cite{efron1979} (400 remuestreos);
la comparación pareada candidato--baseline, con intervalo bootstrap de
la diferencia de macro-F1 y un contraste de permutación de signos
sobre los aciertos por fila (2\,000 permutaciones); y un contraste de
permutación de etiquetas (500 permutaciones, resolución mínima del
p-valor $1/501 \approx 0{,}002$) que verifica que el rendimiento
observado no es compatible con el azar. La selección se declara mejora
significativa únicamente si supera al baseline, el intervalo de la
diferencia pareada excluye el cero y el p-valor del contraste de
signos es inferior a 0{,}05. Cuando varias configuraciones se contrastan
contra el baseline dentro de una misma familia (en la comparación final
del método cuántico, las cinco formulaciones por cada configuración
evaluada), los p-valores pareados se corrigen por multiplicidad con el
procedimiento de Benjamini--Hochberg~\cite{benjamini1995} como ajuste
primario y con la corrección de Holm como control más conservador; un
veredicto de mejora o deterioro exige sobrevivir a esa corrección, no solo
al contraste sin ajustar. Las comparaciones pareadas se calculan sobre
las mismas filas de prueba para todos los métodos; de este modo, la
incorporación posterior del método cuántico no requiere modificar las
configuraciones de referencia y la comparación se mantiene homogénea
fila a fila.

\subsection{Entorno y reproducibilidad}

Todos los experimentos se ejecutan en un entorno Conda con Python 3{,}12
y scikit-learn~\cite{pedregosa2011}, con semillas fijas en todos los
procesos aleatorios y separando la lógica reutilizable de la ejecución
de cada fase. Las condiciones detalladas de reproducción, el coste
computacional por bloque y el índice de trazabilidad que asocia cada
cifra de la memoria con su artefacto de origen se recogen en el
apéndice~\ref{cap:apendice}.
